\subsection{Subpreesquemas lisos de un preesquema liso. Morfismos lisos y morfismos diferencialmente lisos}
\label{subsection:IV.17.12}

\begin{theorem}[17.12.1]
\label{IV.17.12.1}
Sean $f:X\to S$, $h:Y\to S$ dos morfismos localmente de presentación finita, $j:Y\to X$ una inmersión, $y$ un punto de $Y$, $x=j(y)$. Las condiciones siguientes son equivalentes:
\begin{enumerate}
  \item[a)] $h$ es liso en el punto $y$ y $f$ es liso en el punto $x$.
  \item[b)] $f$ es liso en el punto $x$, y el homomorfismo canónico (16.4.21)
  \[
  (\sh{N}_{Y/X})_y\to(j^*(\Omega^1_{X/S}))_y
  \]
  admite un inverso por la izquierda.
  \item[c)] $h$ es liso en el punto $y$ y existe una vecindad abierta $U$ de $y$ en $Y$ tal que $j|U:U\to X$ sea una inmersión regular (16.9.2).

\oldpage[IV-4]{86}

  \item[c')] $h$ es liso en el punto $y$ y existe una vecindad abierta $U$ de $y$ en $Y$ tal que $j|U:U\to X$ sea una inmersión cuasi-regular (es decir (16.9.8), existe una vecindad $U$ de $y$ en $Y$ en la que $\sh{N}_{Y/X}$ es un $\sh{O}_Y$-módulo localmente libre y el homomorfismo canónico
  \[
  \mathbf{S}_{\sh{O}_Y}(\sh{N}_{Y/X})\to\operatorname{gr}^{\bullet}(j)
  \]
  es biyectivo).
\end{enumerate}
\end{theorem}

Para demostrar la equivalencia de a) y b), podemos reducirnos al caso en que $S=\Spec(A)$ y $X=\Spec(B)$ son afines, con $Y=\Spec(B/\mathfrak J)$, donde $\mathfrak J$ es un ideal de tipo finito de $B$, y $B$ es una $A$-álgebra lisa (17.3.2, (ii)). Basta entonces aplicar el criterio jacobiano (0, 22.6.1), así como (0, 19.1.14), teniendo en cuenta que $\Omega^1_{X/S}$ es un $\sh{O}_X$-módulo localmente libre en una vecindad de $x$ y $\sh{N}_{Y/X}$ un $\sh{O}_Y$-módulo de tipo finito (16.1.6).

En segundo lugar, probemos que a) implica c'). Podemos reducirnos nuevamente al caso en que $S$, $X$, $Y$ son afines, $B$ y $B/\mathfrak J$ son $A$-álgebras lisas, y resulta entonces de (17.7.9) y (8.10.5, (iv)) que existen un subanillo noetheriano $A'$ de $A$, una $A'$-álgebra de tipo finito $B'$ y un ideal $\mathfrak J'$ de $B'$ tales que $B=B'\otimes_{A'}A$, $\mathfrak J=\mathfrak J'B$ y que $B'$ y $B'/\mathfrak J'$ son $A'$-álgebras lisas. Observemos que, según (0, 19.3.8), $B'$ sigue siendo una $A'$-álgebra formalmente lisa cuando se dota a $A'$ de la topología discreta y a $B'$ de la topología $\mathfrak J'$-preádica. Por consiguiente (0, 19.5.4), $\mathfrak J'/\mathfrak J'^2$ es un $(B'/\mathfrak J')$-módulo proyectivo y el homomorfismo canónico
\[
\mathbf{S}_{B'/\mathfrak J'}(\mathfrak J'/\mathfrak J'^2)\to\operatorname{gr}_{\mathfrak J'}^{\bullet}(B')
\]
es biyectivo. En otros términos (16.9.8), la inmersión $\Spec(B'/\mathfrak J')\to\Spec(B')$ es cuasi-regular, por tanto regular puesto que $B'$ es noetheriano (16.9.10). Pero, como por hipótesis $B'/\mathfrak J'$ es un $A'$-módulo plano (17.5.1) y $B'$ una $A'$-álgebra de presentación finita, podemos aplicar (11.3.8) reemplazando $\sh{F}$ por $\widetilde{\mathfrak J'}$, y vemos así (por cambio de base) que $j$ es una inmersión regular, y \emph{a fortiori} cuasi-regular. El hecho de que $B$ sea una $A$-álgebra de presentación finita y un $A$-módulo plano muestra además, en virtud de (11.3.8), que las condiciones c') y c) son equivalentes, y son estables por cambio de base.

Demostremos finalmente que c) implica a). Del hecho de que, en (11.3.8), la condición b) implica c), se ve ya que, si c) se cumple, $f$ es \emph{plano} en el punto $x$; en virtud de (17.5.1), todo se reduce, pues, a ver que, bajo la hipótesis c), $f^{-1}(s)$ es liso sobre $k(s)$ en el punto $x$, poniendo $s=h(x)$; como se ha observado que la condición c) es estable por cambio de base, vemos que nos reducimos al caso en que $S=\Spec(k)$ es el espectro de un cuerpo, y teniendo en cuenta (17.7.1, (ii)) podemos suponer que $k$ es \emph{algebraicamente cerrado}. Basta entonces probar que $\sh{O}_{X,x}$ es un anillo \emph{regular} (17.5.1). Ahora bien, por hipótesis se tiene $\sh{O}_{Y,y}=\sh{O}_{X,x}/\mathfrak J_x$, donde $\mathfrak J_x$ es un ideal engendrado por una sucesión $\sh{O}_{X,x}$-regular $(t_i)$, y $\sh{O}_{Y,y}$ es un anillo regular; como los $t_i$ forman parte de un sistema de parámetros para $\sh{O}_{X,x}$ (0, 16.4.1), la conclusión resulta de (0, 17.1.7).

\begin{corollary}[17.12.2]
\label{IV.17.12.2}
Con las notaciones de (17.12.1), supongamos que $f$ es liso en el punto $x$, y sea $Y$ un subpreesquema cerrado de $X$ definido por un ideal $\sh{J}$ de $\sh{O}_X$, y liso sobre $S$ en el punto $x$. Sean $(g_i)_{1\leq i\leq r}$ secciones de $\sh{J}$ sobre $X$. Las condiciones siguientes son equivalentes:
\begin{enumerate}
  \item[a)] Los $(g_i)_x$ forman un sistema de generadores del $\sh{O}_{X,x}$-módulo $\sh{J}_x$ cuyo número de elementos es el menor posible.
  \item[b)] Si $g_i'$ es la imagen canónica de $g_i$ en $\Gamma(X,\sh{J}/\sh{J}^2)$, los $(g_i')_x$ forman una base del $(\sh{O}_{X,x}/\sh{J}_x)$-módulo $\sh{J}_x/\sh{J}_x^2$.

\oldpage[IV-4]{87}

  \item[c)] Las imágenes de los $d_{X/S}g_i$ en $(\Omega^1_{X/S})_x\otimes_{\sh{O}_{X,x}}k(x)$ son elementos linealmente independientes sobre $k(x)$, y los $(g_i)_x$ engendran $\sh{J}_x$.
  \item[d)] Existe una vecindad abierta $U$ de $x$ en $X$ y secciones $g_j$ ($r+1\leq j\leq n$) de $\sh{O}_X$ sobre $U$ tales que el $S$-morfismo $u:U\to S[T_1,\ldots,T_n]$ correspondiente al homomorfismo $\sh{O}_S^n\to f_*(\sh{O}_U)$ definido por las secciones $g_i|U$ ($1\leq i\leq r$) y $g_j$ ($r+1\leq j\leq n$) (II, 1.2.7) sea un morfismo étale, y que la imagen inversa del subpreesquema $Y'=S[T_{r+1},\ldots,T_n]$ por este morfismo sea el preesquema inducido por $j(Y)$ sobre el abierto $U\cap j(Y)$.
\end{enumerate}
\end{corollary}

Como $(g_i')_x$ es la imagen canónica de $(g_i)_x$, la equivalencia de a) y b) resulta del lema de Nakayama, pues $\sh{J}_x$ es de tipo finito y $\sh{J}_x/\sh{J}_x^2$ un $(\sh{O}_{X,x}/\sh{J}_x)$-módulo libre (17.10.4) (Bourbaki, \emph{Alg. comm.}, cap. II, § 3, n\textsuperscript{o} 2, prop. 5). En virtud de (17.12.1, b)), $\sh{J}_x/\sh{J}_x^2$ se identifica canónicamente con un factor directo del $(\sh{O}_{X,x}/\sh{J}_x)$-módulo libre de rango $n$, $(\Omega^1_{X/S})_x\otimes_{\sh{O}_{X,x}}(\sh{O}_{X,x}/\sh{J}_x)$, y la equivalencia de b) y c) resulta de Bourbaki, \emph{loc. cit.} Además, si a) se cumple, $(g_i')_x$ se identifica así con $(d_{X/S}g_i)_x\otimes1$ para $1\leq i\leq r$; como el $\sh{O}_{X,x}$-módulo $(\Omega^1_{X/S})_x$ es libre de rango $n$, resulta nuevamente de Bourbaki, \emph{loc. cit.}, que, reemplazando si es preciso $X$ por una vecindad abierta de $x$, se puede suponer que existen $n-r$ secciones $g_i$ ($r+1\leq i\leq n$) de $\sh{O}_X$ sobre $X$, tales que los $(d_{X/S}g_i)_x$ para $1\leq i\leq n$ forman una base de $(\Omega^1_{X/S})_x$; el hecho de que el morfismo correspondiente $u$ sea étale en el punto $x$ resulta entonces de (17.11.3); por la misma razón, el morfismo $u^{-1}(Y')\to Y'$, restricción de $u$, es étale en el punto $x$; reemplazando $X$ por una vecindad de $x$, se puede suponer que ambos morfismos son étale. Además, es inmediato que $Y$ (identificado, para simplificar, con el subpreesquema cerrado $j(Y)$ de $X$) es un subpreesquema cerrado de $u^{-1}(Y')$, y, en virtud de la elección de los $g_i$ para $i>r$ y de (17.2.5) y (17.11.2), se puede suponer también que la restricción a $Y$ de $u$ es étale. Se deduce que, para todo $y'\in Y'$, la inmersión $Y\cap u^{-1}(y')\to u^{-1}(y')$ es abierta, de donde se concluye, con ayuda de (17.9.6), que $Y\to u^{-1}(Y')$ es una inmersión abierta; esto prueba que a) implica d). El recíproco se deduce inmediatamente de (17.11.3).

\begin{corollary}[17.12.3]
\label{IV.17.12.3}
Sean $f:X\to S$ un morfismo localmente de presentación finita, $u:S\to X$ una $S$-sección de $X$. Para que $f$ sea liso en un punto $x\in X$, es necesario y suficiente que $u$ sea una inmersión cuasi-regular en una vecindad de $s=f(x)$.
\end{corollary}

La conclusión resulta de (17.12.1), puesto que $f\circ u=1_S$ es liso.

\begin{proposition}[17.12.4]
\label{IV.17.12.4}
Todo morfismo liso $f:X\to S$ es diferencialmente liso (es decir (16.10.5), $\Delta_f:X\to X\times_SX$ es una inmersión cuasi-regular). \emph{En particular, los $\sh{P}_f^n=\sh{P}_{X/S}^n$ y los $\operatorname{gr}_n(\sh{P}_{X/S})$ son $\sh{O}_X$-módulos localmente libres de tipo finito.}
\end{proposition}

Basta observar que el morfismo estructural $p_1:X\times_SX\to X$ es liso (17.3.3, (iii)) y que $\Delta_f$ es una $X$-sección de $p_1$; la conclusión resulta de (17.12.3).

\begin{proposition}[17.12.5]
\label{IV.17.12.5}
Sea $f:X\to Y$ un morfismo localmente de presentación finita. Las condiciones siguientes son equivalentes:
\begin{enumerate}
  \item[a)] $f$ es diferencialmente liso.
  \item[b)] Para todo morfismo $g:Y'\to Y$, si se pone $X'=X\times_Y Y'$ y $f'=f_{(Y')}:X'\to Y'$, entonces, para toda $Y'$-sección $s'$ de $X'$, $f'$ es liso en todos los puntos de $s'(Y')$.
  \item[c)] La segunda proyección $p_2:X\times_YX\to X$ es lisa en todos los puntos de la diagonal $\Delta_f(X)$.
\end{enumerate}
\end{proposition}

\oldpage[IV-4]{88}

La condición c) es un caso particular de b): basta, en efecto, tomar $Y'=X$ y $g=f$ en b), pues entonces $f'$ no es sino la segunda proyección $p_2$, y $\Delta_f$ es una $X$-sección de $X\times_YX$. En segundo lugar, c) implica a), pues $p_2$ es un morfismo localmente de presentación finita, y si $p_2$ es liso en los puntos de $\Delta_f(X)$, $\Delta_f$ es una inmersión cuasi-regular (17.12.3), por tanto $f$ es diferencialmente liso. Finalmente, demostremos que a) implica b); si $\Delta_f$ es una inmersión cuasi-regular, $p_2$ es liso en todos los puntos de $\Delta_f(X)$ (17.12.3).

Ahora bien, si $g':X'\to X$ es la proyección canónica, y $v=(g',g')_Y:X'\to X\times_YX$, se comprueba inmediatamente que el diagrama
\[
\xymatrix{
  X\times_Y X \ar[d]_{p_2} & X' \ar[l]_-v \ar[d]^{f'}\\
  X & Y' \ar[l]^{h=g'\circ s'}
}
\]
es conmutativo e identifica $X'$ con el producto $(X\times_YX)\times_XY'$ (I, 3.3.9); teniendo en cuenta que el diagrama (17.4.1.1) identifica $Y'$ con el producto de los $(X\times_YX)$-preesquemas $X$ y $X'$, se concluye, del hecho de que $p_2$ es liso en todo punto de $\Delta_f(X)$, que $f'$ es liso en todo punto de $s'(Y')$ (17.3.3, (iii)).

\begin{corollary}[17.12.6]
\label{IV.17.12.6}
Con las notaciones de (8.8.1), se supone $X_\alpha$ cuasicompacto, y $f_\alpha:X_\alpha\to S_\alpha$ localmente de presentación finita. Para que $f:X\to S$ sea diferencialmente liso, es necesario y suficiente que exista $\lambda\geq\alpha$ tal que $f_\lambda:X_\lambda\to S_\lambda$ sea diferencialmente liso (en cuyo caso $f_\mu:X_\mu\to S_\mu$ es diferencialmente liso para $\mu\geq\lambda$).
\end{corollary}

La suficiencia de la condición y la última afirmación resultan de (16.10.4). Para probar que la condición es necesaria, observemos que $\Delta_f:X\to X\times_SX$ y $p_2:X\times_SX\to X$ se deducen por cambio de base de $\Delta_{f_\lambda}:X_\lambda\to X_\lambda\times_{S_\lambda}X_\lambda$ y $p_{2,\lambda}:X_\lambda\times_{S_\lambda}X_\lambda\to X_\lambda$. En virtud de (17.12.5), para todo $z\in\Delta_f(X)$, existe una vecindad abierta afín $U(z)$ de $z$ en $X\times_SX$ tal que $p_2$ sea liso en $U(z)$; en virtud de (8.2.11) y (17.7.8), existen un índice $\lambda(z)$ y una vecindad $U_{\lambda(z)}$ de la proyección de $z$ en $X_{\lambda(z)}\times_{S_{\lambda(z)}}X_{\lambda(z)}$ tales que $U(z)$ sea la imagen inversa de $U_{\lambda(z)}$, y que $p_{2,\lambda(z)}$ sea liso en $U_{\lambda(z)}$. Como $X$ es cuasicompacto, se puede recubrir $\Delta_f(X)$ por un número finito de vecindades $U(z_i)$; en virtud de (8.3.4), existe un $\lambda$ superior a todos los $\lambda(z_i)$ tal que las imágenes inversas de los $U_{\lambda(z_i)}$ en $X_\lambda\times_{S_\lambda}X_\lambda$ forman un recubrimiento de $\Delta_{f_\lambda}(X_\lambda)$; resulta entonces de (17.12.5) que este índice $\lambda$ cumple lo requerido.

\begin{examples}[17.12.7]
\label{IV.17.12.7}
Sean $S$ un preesquema, $G$ un $S$-\emph{preesquema en grupos}, localmente de presentación finita sobre $S$. Entonces, para que $G$ sea diferencialmente liso, es necesario y suficiente que $G$ sea liso sobre $S$ en los puntos de la «sección unidad» $e$ de $G$. La condición es, en efecto, necesaria por (17.12.5). Recíprocamente, supongamos que se cumple; para todo morfismo $S'\to S$, $G'=G\times_SS'$ es entonces un $S'$-preesquema en grupos localmente de presentación finita sobre $S'$; podemos, por tanto, en virtud de (17.12.5), reducirnos a probar que, para toda $S$-sección $s$ de $G$, $f$ es liso en los puntos de $s(S)$. Ahora bien, si $m:G\times_SG\to G$ es el morfismo que define la estructura de preesquema en grupos de $G$, $m\circ(s\times1_G):S\times_SG\to G\times_SG\to G$ es un

\oldpage[IV-4]{89}

$S$-isomorfismo del preesquema $G$, la «traslación por la izquierda» $\tau_s$, que transforma los puntos de la sección unidad de $G$ en los puntos de $s(S)$; se concluye que $G$ es liso sobre $S$ en los puntos de $s(S)$.
\end{examples}

\begin{remark}[17.12.8]
\label{IV.17.12.8}
Un $S$-preesquema en grupos $G$, de tipo finito (e incluso finito) sobre un preesquema localmente noetheriano $S$, puede ser diferencialmente liso sobre $S$ sin ser liso (ni siquiera \emph{plano}) sobre $S$. Tomemos, por ejemplo, como $S$ el espectro del álgebra de números duales $D=k[T]/T^2k[T]$ sobre un cuerpo $k$, y como $G$ el espectro de la $D$-álgebra producto directo $E=D\oplus k$. Se define en $G$ una estructura de $S$-esquema en grupos definiendo una «aplicación diagonal» $\delta$, homomorfismo de $E$ en $E\otimes_DE$: si $e'$, $e''$ son los idempotentes, imágenes canónicas de la unidad $1$ de $D$ en los factores $D$ y $k$ de $E$, se comprueba fácilmente que se obtiene una aplicación diagonal de este tipo tomando
\[
\delta(e')=e'\otimes e'+e''\otimes e'',\qquad
\delta(e'')=e'\otimes e''+e''\otimes e'.
\]
La sección unidad del esquema en grupos $G$ corresponde al homomorfismo $E\to D$ que es la identidad en $D$ y $0$ en $k$; es un isomorfismo de $S$ sobre una componente conexa de $G$, y \emph{a fortiori} $G$ es diferencialmente liso sobre $S$ (17.12.6), pero está claro que $G$ no es $S$-plano.
\end{remark}
